\part{Strategic Behaviour and Extensions}

\chapter{Game Theory}
\label{chap:game-theory}

\section{Learning objectives}

After completing this chapter, you should be able to:
\begin{itemize}
  \item identify the players, strategies, payoffs, information, and rules in a game;
  \item distinguish an action from a complete strategy;
  \item find dominant strategies, best responses, and Nash equilibria;
  \item analyse sequential games by backward induction and assess whether threats are credible; and
  \item explain how repetition and institutions can change strategic incentives.
\end{itemize}

\section{When choices are interdependent}

\keyterm{Game theory} studies situations in which one person's outcome depends
on their own action \emph{and} on the actions of others. It is useful whenever
people must anticipate one another: firms choosing prices, parties choosing
campaign messages, states deciding whether to negotiate, or students deciding
how much effort to put into a group project.

A game has four basic ingredients:
\begin{itemize}
  \item \keyterm{players}, the decision-makers;
  \item \keyterm{strategies}, the complete plans available to each player;
  \item \keyterm{payoffs}, which rank the consequences for each player; and
  \item rules and information, including who moves when and what each player observes.
\end{itemize}

Payoffs need not be money, and the numbers need not be comparable across
players. They simply represent how each player ranks outcomes in the model.
In a simultaneous game an \keyterm{action} is usually also a complete strategy.
In a sequential game, a strategy must say what a player would do at every
decision point that they might face, including one that is never reached.

\begin{exercisebox}
\textbf{The group-project dilemma.}
Four students value a high-quality group presentation, but preparing it is
costly. Each is tempted to free ride on the others' work. The outcome for one
student therefore depends on several choices, not only on their own effort.
Game theory makes that interdependence explicit before proposing a rule, a
grade scheme, or a norm to improve the result.
\end{exercisebox}

\begin{keyidea}
\textbf{A model of strategic interaction is not a moral judgement.}
Calling an action a best response says that it is privately advantageous under
the stated rules and beliefs. It does not say that the action is fair, socially
desirable, or inevitable once the rules can be changed.
\end{keyidea}

\section{Strategic-form games}

A \keyterm{strategic-form game} represents simultaneous choices in a payoff
matrix. Rows record one player's strategies, columns record the other's, and
each ordered pair gives the payoffs in the order \((\text{row player},
\text{column player})\).

\subsection{Best responses, dominance, and Nash equilibrium}

A \keyterm{best response} is a strategy that gives a player the highest payoff
given the other players' strategies. A \keyterm{strictly dominant strategy}
gives a higher payoff than every alternative for every possible strategy of the
opponent. A \keyterm{strictly dominated strategy} is always worse than another
available strategy, so a rational player would never need to use it.

A \keyterm{Nash equilibrium} is a profile of strategies in which every player
is choosing a best response to the others. Formally, a profile
\(s^*=(s_i^*,s_{-i}^*)\) is a Nash equilibrium when
\[
u_i(s_i^*,s_{-i}^*)\geq u_i(s_i,s_{-i}^*)
\qquad\text{for every feasible }s_i.
\]
Here \(s_{-i}^*\) means the strategies chosen by everyone other than player
\(i\). At a Nash equilibrium, no individual can improve their own payoff by
changing strategy alone. The concept predicts stability; it does \emph{not}
guarantee cooperation, fairness, or efficiency.

\begin{keyidea}
\textbf{How to read a payoff matrix.}
\begin{enumerate}
  \item Hold the column fixed and compare the first number: this finds Player~1's best responses.
  \item Hold the row fixed and compare the second number: this finds Player~2's best responses.
  \item A cell in which both choices are best responses is a Nash equilibrium.
\end{enumerate}
\end{keyidea}

\subsection{The Prisoner's Dilemma}

The \keyterm{Prisoner's Dilemma} is the classic case in which individually
rational choices lead to a jointly inferior outcome. In Figure~\ref{fig:prisoners-dilemma},
Defect is strictly dominant for both players. The gold-bordered cell is
therefore the unique Nash equilibrium. Yet both players prefer mutual
cooperation to mutual defection.

\begin{figure}[htbp]
  \centering
  \begin{tikzpicture}[font=\small]
    \node[font=\sffamily\bfseries] at (4.0,3.05) {Player 2};
    \node[font=\sffamily] at (2.9,2.55) {Cooperate};
    \node[font=\sffamily] at (5.1,2.55) {Defect};
    \node[font=\sffamily\bfseries,rotate=90] at (0.45,1.2) {Player 1};
    \node[font=\sffamily,anchor=east] at (1.65,1.8) {Cooperate};
    \node[font=\sffamily,anchor=east] at (1.65,0.6) {Defect};
    \draw[PPEInk] (1.8,0) rectangle (6.2,2.4);
    \draw[PPEInk] (4.0,0) -- (4.0,2.4);
    \draw[PPEInk] (1.8,1.2) -- (6.2,1.2);
    \node at (2.9,1.8) {\colorbox{PPEAccentLight}{\strut $(3,3)$}};
    \node at (5.1,1.8) {$(0,5)$};
    \node at (2.9,0.6) {$(5,0)$};
    \node at (5.1,0.6) {\fcolorbox{PPEGold}{PPEGold!18}{\strut $(1,1)$}};
    \node[font=\scriptsize,text=PPEGold] at (5.1,-0.43) {Nash equilibrium};
    \node[font=\scriptsize,text=PPEBurgundy,align=center] at (2.9,-0.56)
      {mutually preferred\\cooperation};
  \end{tikzpicture}
  \caption{A Prisoner's Dilemma. Payoffs are ordered
  \((\text{Player 1},\text{Player 2})\). Mutual defection is a Nash
  equilibrium even though mutual cooperation Pareto dominates it.}
  \label{fig:prisoners-dilemma}
\end{figure}
\FloatBarrier

This logic appears in cartel formation, climate policy, arms races, tax
compliance, and team production. It is not enough to tell players that
cooperation would be better for both. Cooperation must be supported by
contracts, monitoring, repeated interaction, changed payoffs, or a shared norm.

\section{Coordination and multiple equilibria}

Not every strategic problem is a conflict. In a \keyterm{coordination game},
players benefit from selecting compatible actions but may disagree about which
compatible outcome is best. The \keyterm{Battle of the Sexes} is a compact
example:
\[
\begin{array}{c|cc}
 & \text{Player 2: Opera} & \text{Player 2: Football}\\
\hline
\text{Player 1: Opera} & (2,1) & (0,0)\\
\text{Player 1: Football} & (0,0) & (1,2)
\end{array}
\]
There are two pure-strategy Nash equilibria: both choose Opera and both choose
Football. Neither player has a dominant strategy. A prior agreement, a social
convention, a visible focal point, or a rule that allocates turns can therefore
create value by helping players coordinate.

The possibility of multiple equilibria matters for public policy. A rule can
matter not only by changing payoffs, but also by making one expectation more
salient or by supplying a reliable way to coordinate.

\subsection{More games}
\label{sec:more-games}

The exercise set for this course contains several useful variants of the
games introduced so far. Their names are less important than their incentive
structures. The following examples add five distinct structures: risky
coordination, anti-coordination, zero-sum competition, strategic pricing with
more than two actions, and a larger political game.

\subsubsection{Stag Hunt: coordination with an assurance problem}

In a \keyterm{Stag Hunt}, both players prefer the high-payoff outcome when
they coordinate on it, but each has a safer action if they doubt the other
player's commitment:
\[
\begin{array}{c|cc}
 & \text{Player 2: Stag} & \text{Player 2: Hare}\\
\hline
\text{Player 1: Stag} & (10,10) & (0,8)\\
\text{Player 1: Hare} & (8,0) & (5,5)
\end{array}
\]
There are two pure-strategy Nash equilibria: \((\text{Stag},\text{Stag})\)
and \((\text{Hare},\text{Hare})\). Stag is collectively more productive, but
Hare protects a player against being the only one to choose Stag. If Player~1
believes that Player~2 chooses Stag with probability \(p\), Stag gives
\(10p\), whereas Hare gives \(8p+5(1-p)\). Thus Player~1 chooses Stag only
when \(p>5/7\). The central obstacle is therefore assurance: each player needs
reason to expect the other to make the ambitious choice. Investment,
technology adoption, and collective action often have this form.

\subsubsection{Chicken: anti-coordination and brinkmanship}

\keyterm{Chicken}, also called a Hawk--Dove game in many applications, has
the opposite logic. Each player prefers to take the tough action when the
other gives way, but both are harmed if neither gives way:
\[
\begin{array}{c|cc}
 & \text{Player 2: Stop} & \text{Player 2: Continue}\\
\hline
\text{Player 1: Stop} & (0,0) & (-5,10)\\
\text{Player 1: Continue} & (10,-5) & (-15,-15)
\end{array}
\]
The two pure-strategy Nash equilibria are
\((\text{Stop},\text{Continue})\) and
\((\text{Continue},\text{Stop})\). Unlike a coordination game, the players
want to choose opposite actions. The mutual-Continue outcome is especially
bad, so a credible commitment, a turn-taking rule, or communication can avoid
a costly collision. The game is a stylised way to analyse bargaining,
brinkmanship, or rival firms' aggressive campaigns; it is not a prediction
that real conflict must be dangerous.

\subsubsection{Rock--Paper--Scissors: zero-sum competition and mixing}

In a \keyterm{zero-sum game}, one player's gain is the other's loss. The
course exercise on Rock--Paper--Scissors has the following payoffs, where
\(\mathrm{R}\), \(\mathrm{S}\), and \(\mathrm{P}\) stand for Rock, Scissors,
and Paper:
\[
\begin{array}{c|ccc}
 & \text{P2: }\mathrm{R} & \text{P2: }\mathrm{S}
 & \text{P2: }\mathrm{P}\\
\hline
\text{P1: }\mathrm{R} & (0,0) & (10,-10) & (-10,10)\\
\text{P1: }\mathrm{S} & (-10,10) & (0,0) & (10,-10)\\
\text{P1: }\mathrm{P} & (10,-10) & (-10,10) & (0,0)
\end{array}
\]
There is no pure-strategy Nash equilibrium: after every predictable choice,
the other player can switch to the winning response. The equilibrium is a
\keyterm{mixed strategy}: each player deliberately chooses Rock, Scissors,
and Paper with probability \(1/3\). Randomisation here is not indecision. It
prevents an opponent from exploiting a predictable pattern. A penalty
shoot-out is a two-action version of the same idea: in the symmetric exercise,
the kicker and goalkeeper each choose Left and Right with probability \(1/2\).

\subsubsection{Bar prices: iterated elimination in a \(3\times3\) game}

Strategic-form games need not have only two actions. In the bar-price exercise,
La Piedra chooses a row and Shamrock chooses a column. The prices are in euros
and the ordered pairs are illustrative profit indices:
\[
\begin{array}{c|ccc}
 & \multicolumn{3}{c}{\text{Shamrock's price}}\\
\text{La Piedra's price} & 2\,\euro & 4\,\euro & 5\,\euro\\
\hline
2\,\euro & (10,10) & (14,12) & (14,15)\\
4\,\euro & (12,14) & (20,20) & (28,15)\\
5\,\euro & (15,14) & (15,28) & (25,25)
\end{array}
\]
For each bar, \(2\,\euro\) is strictly dominated by \(4\,\euro\), so both
players can remove that option. In the reduced game, \(4\,\euro\) strictly
dominates \(5\,\euro\) for both players. This
\keyterm{iterated elimination of strictly dominated strategies} therefore
leads to the unique Nash equilibrium \((4\,\euro,4\,\euro)\). The example
shows why a price that looks attractive in isolation need not be a best
response to a competitor's price.

\subsubsection{Arms race: dominance in a political game}

The exercise set also uses a deliberately simplified USA--USSR arms-race
matrix. Each player can Demilitarise (\(\mathrm{D}\)), Escalate
(\(\mathrm{E}\)), or Go Nuclear (\(\mathrm{GN}\)):
\[
\begin{array}{c|ccc}
 & \text{USSR: \(\mathrm{D}\)} & \text{USSR: \(\mathrm{E}\)}
 & \text{USSR: \(\mathrm{GN}\)}\\
\hline
\text{USA: \(\mathrm{D}\)} & (20,30) & (10,40) & (-10,50)\\
\text{USA: \(\mathrm{E}\)} & (40,-10) & (10,20) & (-5,15)\\
\text{USA: \(\mathrm{GN}\)} & (50,-100) & (30,-10) & (5,0)
\end{array}
\]
For the USSR, \(\mathrm{D}\) is strictly dominated by
\(\mathrm{E}\). Once \(\mathrm{D}\) is removed, \(\mathrm{GN}\) strictly
dominates the USA's remaining alternatives. The USSR's best response to USA
\(\mathrm{GN}\) is then \(\mathrm{GN}\), giving the unique equilibrium
\((\mathrm{GN},\mathrm{GN})\). The conclusion follows from the stated
payoffs, not from a historical claim about the Cold War. A PPE analysis should
ask who chose these preference rankings, which risks or ethical costs they
omit, and how treaties or verification could change the game.

\begin{keyidea}
\textbf{Names do not solve games; incentives do.}
The marketing-campaign exercise has the same dominant-strategy structure as
the Prisoner's Dilemma already analysed in this chapter. The plain
coordination matrix and Battle of the Sexes are also already represented
above, while the entry game and the Centipede game introduce the sequential
cases. For any new matrix, use the same method: identify dominated strategies,
find best responses, locate equilibria, and then interpret the assumptions
behind the payoffs.
\end{keyidea}

\section{Sequential games, credibility, and backward induction}

In a \keyterm{dynamic game}, players move in sequence. An
\keyterm{extensive-form game} uses a tree to show the order of moves,
information, available actions, and terminal payoffs. \keyterm{Backward
induction} solves a finite extensive-form game with perfect information from
the final decision nodes back to the first move.

Consider entry into a market. If the entrant stays out, payoffs are
\((0,3)\): the entrant earns zero and the incumbent retains monopoly profit.
If the entrant enters, the incumbent can Fight, yielding \((-1,-2)\), or
Accommodate, yielding \((1,1)\). Figure~\ref{fig:entry-deterrence} shows why a
threat to fight is not credible: once entry has occurred, accommodation gives
the incumbent \(1\), rather than \(-2\).

\begin{figure}[htbp]
  \centering
  \begin{tikzpicture}[x=1cm,y=1cm,font=\small]
    \fill[PPEInk] (0,0) circle (1.6pt);
    \node[font=\sffamily\bfseries,above] at (0,0) {Entrant};
    \draw[->,PPEInk] (0.11,0.04) -- (2.75,1.02)
      node[midway,above,font=\scriptsize] {Out};
    \node[anchor=west] at (2.9,1.02) {$(0,3)$};
    \draw[->,PPENavarraRed,very thick] (0.11,-0.04) -- (2.75,-1.02)
      node[midway,below,font=\scriptsize] {Enter};
    \fill[PPEInk] (2.85,-1.06) circle (1.6pt);
    \node[font=\sffamily\bfseries,below] at (2.85,-1.06) {Incumbent};
    \draw[->,PPEInk] (2.98,-1.11) -- (5.65,-2.06)
      node[midway,below,font=\scriptsize] {Fight};
    \node[anchor=west] at (5.8,-2.06) {$(-1,-2)$};
    \draw[->,PPENavarraRed,very thick] (2.98,-1.01) -- (5.65,-0.16)
      node[midway,above,font=\scriptsize] {Accommodate};
    \node[anchor=west] at (5.8,-0.16) {$(1,1)$};
    \node[font=\scriptsize,text=PPEBurgundy,align=center] at (4.9,-2.9)
      {subgame-perfect equilibrium:\\Enter, Accommodate};
  \end{tikzpicture}
  \caption{Entry deterrence. The highlighted path follows backward induction:
  after entry, the incumbent prefers accommodation, so the entrant enters.}
  \label{fig:entry-deterrence}
\end{figure}
\FloatBarrier

A \keyterm{credible threat} is an action that is optimal when the relevant
contingency arrives. A \keyterm{subgame-perfect equilibrium} is a strategy
profile that forms a Nash equilibrium in every subgame. It rules out an empty
threat of fighting after entry when fighting would actually harm the incumbent.

\section{Repeated interaction}

Repeated play can change the incentives in a one-shot Prisoner's Dilemma. If
players expect to meet again, a short-run gain from defection can be outweighed
by the loss of future cooperation. This is one reason why reputation,
long-lived relationships, and transparent records can matter in markets.

For a simple illustration, suppose each player receives \(3\) every period if
both cooperate. A player who defects while the other cooperates receives \(5\)
once, but then receives \(1\) forever after. Let \(\delta\) be the
\keyterm{shadow of the future}: the value placed on a payoff one period ahead
relative to a payoff today. Cooperation is sustainable under this simple
trigger strategy if
\[
\frac{3}{1-\delta}\geq 5+\frac{\delta}{1-\delta},
\qquad\text{which implies}\qquad
\delta\geq \frac{1}{2}.
\]
The calculation is not a universal law. It assumes reliable observation,
credible punishment, an indefinite or uncertain ending date, and patient
players. In a known final period, backward induction can instead unravel
cooperation.

The analysis connects directly to Chapter~13. Competing firms may all gain
from keeping prices high, but each has an incentive to cut price or expand
output. Competition law, changing market conditions, entry, and imperfect
monitoring can all make collusion difficult to sustain.

\section{Institutions change games}

Game theory helps compare institutions because a rule can change a game in
several ways: it can alter payoffs through taxes or penalties, reveal
information through disclosure, change which actions are feasible through
contracts, or coordinate expectations through a clear procedure. The central
PPE question is therefore not whether strategic behaviour is good or bad, but
which institutional design produces the most defensible outcomes.

\section*{Chapter Review}

\subsection*{Key concepts and equations}
\begin{itemize}
  \item \textbf{Strategy}: a complete contingent plan available to a player.
  \item \textbf{Best response}: a payoff-maximising choice given the other players' choices.
  \item \textbf{Nash equilibrium}: a profile of mutual best responses.
  \item \textbf{Prisoner's Dilemma}: individually rational defection can produce an inefficient outcome.
  \item \textbf{Coordination game}: players benefit from compatible choices and may face multiple equilibria.
  \item \textbf{Backward induction}: solving a finite perfect-information
  game from final decision nodes back to earlier choices.
  \item \textbf{Credible threat}: an action that remains optimal when the contingency occurs.
  \item \textbf{Subgame-perfect equilibrium}: a Nash equilibrium in every subgame.
  \item \textbf{Repeated game}: strategic interaction that occurs more than once.
\end{itemize}

\section*{Exercises}

\subsection*{Conceptual questions}
\begin{enumerate}[label=\arabic*.]
  \item Identify the players, strategies, payoffs, and rules in one economic or political interaction.
  \item Distinguish an action from a strategy in a sequential game.
  \item Distinguish a dominant strategy from a best response.
  \item Define a Nash equilibrium. Why need it not be Pareto efficient?
  \item Explain why a threat can fail to be credible.
  \item How can repeated interaction support cooperation, and what conditions can prevent it?
\end{enumerate}

\subsection*{Graphical exercises}
\begin{enumerate}[resume,label=\arabic*.]
  \item Construct a \(2\times2\) Prisoner's Dilemma payoff matrix. Mark each player's best responses, its Nash equilibrium, and an outcome that Pareto dominates it.
  \item Draw the entry-deterrence game in Figure~\ref{fig:entry-deterrence}. Change the payoffs so that Fight becomes credible, then solve the revised game by backward induction.
\end{enumerate}

\subsection*{Analytical exercises}
\begin{enumerate}[resume,label=\arabic*.]
  \item In the Battle of the Sexes matrix above, identify all pure-strategy Nash equilibria. Does either player have a dominant strategy?
  \item An entrant chooses Enter or Stay Out. If Enter, an incumbent chooses Fight or Accommodate. The payoffs are Out \((0,3)\), Enter--Fight \((-1,-2)\), and Enter--Accommodate \((1,1)\). Find the subgame-perfect equilibrium.
  \item In the repeated Prisoner's Dilemma in this chapter, derive the minimum value of \(\delta\) that makes cooperation sustainable under the trigger strategy.
  \item Explain how the temptation to deviate from a cartel agreement resembles a Prisoner's Dilemma. Give one reason why a real cartel may be harder to sustain than the simple repeated-game model suggests.
\end{enumerate}

\subsection*{Exam-style multiple-choice questions}
\begin{enumerate}[resume,label=\arabic*.]
  \item A Nash equilibrium is an outcome where: (a) every player cooperates; (b) no player benefits by unilaterally changing action; (c) all payoffs are equal; or (d) one player wins.
  \item In a finite dynamic game, backward induction begins by analysing: (a) the first move; (b) the final decision nodes; (c) average payoffs; or (d) market demand.
  \item A subgame-perfect equilibrium excludes: (a) all cooperation; (b) credible actions; (c) non-credible threats; or (d) sequential decisions.
\end{enumerate}

\subsection*{PPE discussion question}

\noindent Choose a strategic problem in policy, business, or politics. Which
aspect of the game should an institution change: payoffs, information, feasible
strategies, or expectations? Defend one realistic proposal and discuss a
possible unintended consequence.

\subsection*{Optional extension: the Centipede game}

\begin{exercisebox}
\textbf{Why do people sometimes pass?}
In a \keyterm{Centipede game}, two players take turns either to Take a current
payoff or to Pass, which enlarges the available payoff for the next decision.
In a finite version, backward induction says that the last player takes; the
same logic then works back to the first decision. Yet in laboratory versions,
many people pass early. The contrast is instructive: when a prediction and
observed behaviour differ, ask whether the horizon, beliefs, social
preferences, or interpretation of payoffs matches the model.

\smallskip
\textit{Try it:} construct a four-move Centipede game in which the total
available payoff grows after every Pass. Solve it by backward induction, then
explain one reason why two real people might nevertheless pass at the first
move.
\end{exercisebox}
